% Kapitel 1 — Budget: 6 Seiten — Schreiben: 14.–23.12.
% Quelle Gliederung: codebase-graph/docs/notizen/"Zeitplan und Gliederung 2026-10-19.md" (Stand 26.10.)
% Text für 1.1–1.5 liegt im Exposé vom 21.08.2026 vor und wird überarbeitet übernommen.

\chapter{Einleitung}
\label{chap:einleitung}

\section{Problemstellung und Motivation}% 1 S., aus Exposé überarbeiten
\label{sec:problemstellung}

\section{Das Reaktivitätsproblem}% 1 S., aus Exposé überarbeiten
\label{sec:reaktivitaetsproblem}

\section{Das Strukturproblem}% 1 S., aus Exposé überarbeiten
\label{sec:strukturproblem}

\section{Das Transparenzproblem}% 1 S., aus Exposé überarbeiten
\label{sec:transparenzproblem}

\section{Zielsetzung und Forschungsfragen}% 1,5 S.
% Inhalt: FF1 bis FF6 wie im Exposé; der Wechsel zur Hauptfrage der Studie
% (Zusammenhänge verstehen, Änderung nachvollziehen; Wortlaut 6.1, 13.11.)
% wird in 5.6 und 6.1 begründet.
\label{sec:forschungsfragen}
% FF1 mit E1 bis E3 am 04.11. wörtlich aus dem Exposé übernommen
% (Expose_v4/sections/02_forschungsfragen.tex, Unterabschnitt "Gestaltung des Artefakts").
% Noch nicht überarbeitet: Semikolons und Gedankenstriche aus dem Exposé stehen noch drin.
% 4.2 (sec:konzeption) greift E1 bis E3 in Kurzform auf und verweist hierher.
% =====================================================================
% PFLICHT (13.11.) -- MUSS IN 1.5 STEHEN
% 6.1 (sec:studienfragen) verweist für die neue Hauptfrage der Studie auf diesen Abschnitt und braucht hier drei Dinge im Text, sonst zeigt der Verweis ins Leere:
% 1. Zielsetzung mit dem Gedanken aus dem Exposé (Expose_v4/sections/01_problemstellung.tex, Unterabschnitt "Zielsetzung", zweiter Absatz): Der Entwickler gestaltet die Auswahl, bevor der Agent damit weiterarbeitet, und "muss den Ablauf nicht erst im Nachhinein nachvollziehen". Darauf stützt sich der zweite Teil der Hauptfrage (Änderung des Agenten nachvollziehen).
% 2. Leitfrage des Exposés (02_forschungsfragen.tex, Anfang) mit dem Vergleich kuratiert gegen rein automatisch. 6.1 sagt, die Hauptfrage übernehme den Vergleich zwischen kuratiertem und automatisch abgerufenem Kontext "aus FF3 und FF4". 6.1 verweist außerdem für FF2 bis FF6 hierher (tab:ff-zuordnung), im Text stehen bisher nur FF1 und E1 bis E3.
% 3. Bei FF6 den Absatz zum Nebennutzen mitnehmen ("setzt sich dabei mit der Codebase auseinander und kennt sie hinterher womöglich besser"). Darauf stützt sich der erste Teil der Hauptfrage, und 6.1 zitiert den Satz fast wörtlich.
% 4. FF2 bis FF6 im Wortlaut des Exposés. tab:ff-zuordnung in 6.1 und tab:messgroessen in 6.4 nennen FF2 bis FF6 und verweisen hierher.
% ENDE PFLICHT (13.11.)
% =====================================================================

Das Transparenzproblem verlangt, dass der Entwickler nicht nur das Ergebnis des Abrufs sieht, sondern dessen Zustandekommen. Daraus folgt:

\begin{quote}
\textbf{FF1:} Wie ist ein graphbasiertes Retrieval zu gestalten, das zu jedem abgerufenen Knoten den Auswahlgrund offenlegt und die Auswahl vor der Übergabe an den Agenten schrittweise modifizierbar macht?
\end{quote}

Als Auswahlgrund gelten der Ähnlichkeitswert des Knotens, die wirksame Schnittgrenze, im Ablehnungsfall der Grund der Verwerfung, der eingebettete Beschreibungstext sowie der Traversierungspfad, über den der Knoten erreicht wurde.

Konstruktionsfragen werden nicht über einen Bedingungsvergleich beantwortet. Das gestaltungsorientierte Vorgehen verlangt jedoch, dass ein Artefakt an vorab benannten Entwurfszielen geprüft wird \parencite{hevner2004dsr}. Ohne solche Ziele bliebe FF1 eine Anforderungsliste ohne Kriterium für ihre Erfüllung. Für das hier entwickelte System gelten drei:

\begin{description}[leftmargin=2.2em,style=nextline,itemsep=0.2em]
  \item[E1] Zu jedem Knoten des übergebenen Subgraphen ist der Auswahlgrund im oben genannten Umfang abrufbar, einschließlich der Gründe für verworfene Knoten.
  \item[E2] Jede Entscheidung des Abrufs --- Einstiegspunkte, verfolgte Kanten, Traversierungstiefe, Auswahl und Gewichtung --- ist vor der Übergabe an den Agenten veränderbar.
  \item[E3] Die offengelegten Begründungen werden bei der Kuratierung tatsächlich herangezogen.
\end{description}

E1 und E2 sind Eigenschaften der Konstruktion und werden durch die Beschreibung des Systems belegt. E3 ist es nicht: Ob die Offenlegung genutzt wird oder lediglich vorhanden ist, entscheidet sich am Verhalten der Entwickler. Dieses Ziel wird daher empirisch geprüft, über das Kuratierungsprotokoll, das laute Denken und den Interviewleitfaden der Nutzerstudie. Werden Ähnlichkeitswert und Traversierungspfad durchgängig übergangen, ist E3 verfehlt --- ein Befund, der für die Gestaltung solcher Werkzeuge aufschlussreicher wäre als seine Bestätigung.

\section{Aufbau der Arbeit}% 0,5 S.
% Inhalt: Rollen von Claude (Entwicklungswerkzeug, gemessener Agent, Schreibhilfe)
% und Kennzeichnung der KI-Nutzung nach HTWG-Regel.
\label{sec:aufbau}
